\documentclass[12pt]{article}

\usepackage[a4paper, margin=2.5cm]{geometry}
\usepackage{setspace}
\usepackage{times}
\usepackage{amsmath}
\usepackage{amssymb}
\usepackage{hyperref}
\usepackage{microtype}
\usepackage{parskip}
\usepackage{titlesec}
\usepackage{enumitem}

\singlespacing

\titleformat{\section}{\normalfont\large\bfseries}{\thesection.}{0.5em}{}
\titleformat{\subsection}{\normalfont\normalsize\bfseries}{\thesubsection}{0.5em}{}

\hypersetup{colorlinks=true, linkcolor=black, citecolor=black, urlcolor=black}

% Manual citation commands
\newcommand{\citep}[1]{\textup{(\citeauthor{#1})}}

% Define references manually
\usepackage{etoolbox}

\begin{document}

% Title -- double-blind version
\begin{titlepage}
\begin{center}
\vspace*{2cm}
{\LARGE\bfseries Irreversibility-Constrained Governance}\\[0.8em]
{\large Preserving Human Responsibility under Finite Oversight Capacity}\\[3cm]
{\normalsize \today}
\end{center}
\end{titlepage}

% Abstract
\begin{abstract}
AI-mediated governance environments are characterized by continuous updating and adaptive
optimization. While such systems enhance predictive performance, they introduce a structural
risk: the progressive erosion of human oversight capacity prior to irreversible outcomes.

Whereas frameworks addressing temporal fixation of responsibility specify when institutional
commitment crystallizes, Irreversibility-Constrained Governance (ICG) specifies how institutions
regulate the trajectory toward that fixation under conditions of finite human capacity.

ICG models institutional irreversibility as a function of remaining oversight capacity and
constrains the rate at which systems approach configurations in which evaluative and corrective
judgment collapses. When predefined thresholds are exceeded, a structured freeze mechanism
enforces recovery conditions that restore oversight viability.

The framework is instantiated in a high-risk clinical decision support system for sepsis
management, illustrating how update regulation and recovery operations preserve the evaluative
space necessary for responsible clinical judgment.

The model does not optimize growth. It constrains approach velocity toward institutional
irreversibility. Its objective is to preserve the structural conditions under which human
responsibility remains operational rather than merely formal.

\noindent\textbf{Keywords:} AI governance; institutional irreversibility; oversight capacity;
medical AI; velocity constraint; freeze mechanism; collapse-delay model; human responsibility
\end{abstract}

\newpage

\section{Introduction}

Artificial intelligence systems deployed in institutional environments increasingly operate
under conditions of continuous retraining, parameter recalibration, and adaptive optimization.
These developments enhance predictive accuracy and responsiveness. They also accelerate system
evolution.

Existing responsibility frameworks emphasize that meaningful human control must be ensured
and that system outputs remain responsive to relevant human reasons
(Santoni de Sio \& van den Hoven, 2018).
Such accounts clarify alignment and attribution conditions. However, they do not model how
oversight capacity itself evolves over time in adaptive environments.

A companion framework addressing the temporal fixation of institutional responsibility under
AI-mediated conditions has been developed in parallel work. The present analysis operates
at a prior analytical layer.

Irreversibility-Constrained Governance (ICG) addresses the dynamic conditions under which
institutional oversight capacity may deplete before irreversible outcomes occur.

The central claim is structural: governance must regulate trajectories, not merely outcomes.
A system that approaches irreversibility faster than oversight capacity can be restored risks
rendering responsibility formally assigned yet substantively hollow.

ICG does not oppose optimization. It constrains the velocity of approach toward institutional
irreversibility in order to preserve the institutional ability to stop.

\section{Two Orders of Irreversibility}

\subsection{Biological Irreversibility}

Biological irreversibility refers to outcomes that cannot be undone at the level of
physiological reality, including death, permanent organ damage, and irreversible neurological
impairment. These represent absolute thresholds beyond which corrective intervention
loses restorative meaning.

\subsection{Institutional Irreversibility}

Institutional irreversibility refers to configurations in which evaluative and supervisory
capacity degrade to such an extent that corrective intervention becomes structurally
ineffective. Indicators include:

\begin{itemize}[leftmargin=1.5em]
  \item Persistent update opacity
  \item Review overload exceeding institutional tolerance
  \item Responsibility diffusion across actors and time
  \item Acceleration of system evolution beyond interpretive capacity
  \item Structural inability to suspend or recalibrate system behavior
\end{itemize}

Institutional irreversibility is not harm itself. It is the erosion of the structural
conditions required to prevent harm.

\subsection{Structural Relation}

Institutional irreversibility increases the probability of biological irreversibility by
reducing the evaluative space necessary for timely intervention. ICG therefore targets
the approach to institutional irreversibility. By preserving oversight capacity, institutions
preserve the structural possibility of preventing irreversible harm.

\section{Formal Model of Institutional Irreversibility}

\subsection{Oversight Capacity}

Let institutional oversight capacity be defined as:
\[
B(t) \in [0,1]
\]
where $B(t) = 1$ denotes full evaluative and corrective capacity and $B(t) = 0$ denotes
structural collapse of effective oversight. Oversight capacity aggregates attention,
interpretive bandwidth, audit resources, and intervention authority.

\subsection{Institutional Irreversibility Distance}
\[
D_{II}(t) := 1 - B(t)
\]
This metric captures proximity to institutional configurations in which corrective
intervention becomes structurally ineffective.

\subsection{Capacity Dynamics}
\[
\frac{dB}{dt} = r(t) - c(t)
\]
where $c(t) \geq 0$ is the consumption rate generated by alerts, updates, and review burden;
and $r(t) \geq 0$ is the recovery rate generated by structured institutional interventions.

The dynamics of $B(t)$ are not symmetric. Consumption $c(t)$ may increase rapidly in
response to system updates, alert spikes, or exception cascades, whereas recovery $r(t)$
is constrained by institutional factors such as staff availability, audit backlog, and
procedural latency.

This asymmetry introduces the possibility of institutional hysteresis. When oversight
capacity depletes abruptly, restoration to prior numerical levels of $B(t)$ does not
necessarily restore the evaluative context within which prior judgments were formed.
Recovery may require time exceeding the governance window within which intervention
remains practically meaningful.

Moreover, cumulative drift can compound this effect. Repeated cycles of depletion and
partial recovery may leave the institution structurally closer to irreversibility than
instantaneous values of $B(t)$ alone would indicate.

For this reason, ICG treats the velocity constraint as primary. Preventing rapid depletion
is institutionally more robust than relying on recovery after substantial capacity loss.
Accordingly, $v_{\max}$ must be set conservatively, and freeze activation precedes collapse
rather than reacting to it.

\subsection{Velocity Constraint}
\[
\frac{dD_{II}}{dt} \leq v_{\max}
\]
Equivalently:
\[
c(t) - r(t) \leq v_{\max}
\]
Institutions must not approach irreversibility faster than they can supervise.
This is a governance constraint, not a claim of mechanical stability.

\subsection*{3.4a\quad Governance Interpretation of the Velocity Constraint}

The velocity constraint formalizes an institutional requirement: the rate at which oversight
capacity depletes must not exceed the rate at which institutions can meaningfully respond.

$v_{\max}$ is not derived from engineering specifications. It is a governance parameter set
institutionally based on staff capacity, review infrastructure, and decision complexity.
In medical environments, where biological irreversibility represents an absolute downstream
threshold, $v_{\max}$ must be set conservatively.

The governing variable is not accuracy or efficiency. It is the remaining margin within
which responsible judgment remains structurally possible.

\subsection{Structured Freeze}

If either $D_{II}(t) \geq \theta$ or $\frac{dD_{II}}{dt} \geq \kappa$, structured freeze
activates. Freeze operations:

\begin{itemize}[leftmargin=1.5em]
  \item Suspend model updates
  \item Limit alert frequency
  \item Redistribute review responsibility
  \item Enforce human-only fallback
  \item Restore audit visibility
\end{itemize}

Freeze enforces $r(t) \geq r_{\min}$, restoring the velocity constraint.
Freeze is not prohibition. It is recovery enforcement.

\section{Clinical Instantiation: ICU Sepsis Decision Support}

Deployment of machine learning systems in ICU environments has been shown to impose
significant interpretive and supervisory burden on clinical staff (Sendak et al., 2020).

Regulatory guidance for adaptive medical AI, including the FDA's SaMD Action Plan
(FDA, 2021a), outlines structured change-control processes. Good Machine Learning
Practice (FDA, 2021b) articulates implementation principles for validation and oversight.
However, these frameworks do not model oversight capacity as a dynamically depleting variable.

\subsection{Operationalizing Capacity}
\[
B(t) = 1 - \frac{L(t)}{L_{\max}}
\]
where $L(t)$ denotes review load in minutes per hour and $L_{\max}$ denotes institutional
maximum review tolerance. Institutional irreversibility distance becomes:
\[
D_{II}(t) = \frac{L(t)}{L_{\max}}
\]

\subsection{Consumption Components}
\[
c(t) = \alpha A(t) + \beta U(t) + \gamma E(t)
\]
with alert frequency $A(t)$, update frequency $U(t)$, and exception reviews $E(t)$.

\subsection{Threshold and Freeze}

Let $\theta = 0.6$. If review burden exceeds 60\% of institutional tolerance: model
updates suspend, alert triage activates, human-only fallback is enforced, and audit
locking is initiated. ICG regulates the structural conditions of judgment, not clinical
outcomes directly.

\subsection{Clinical Scenario: Capacity Depletion and Recovery}

Consider a twelve-hour ICU shift in which a sepsis decision support system undergoes
two unscheduled model updates following new patient cohort data. Each update alters
risk threshold distributions, requiring clinical staff to re-evaluate prior recommendations.

At $t = 0$, $B(t) = 0.85$. Review load is within tolerance.

Between $t = 1$ and $t = 3$, alert frequency $A(t)$ increases as the updated model flags
a broader patient population. Consumption rate $c(t)$ rises. Recovery rate $r(t)$ remains
constant. $D_{II}(t)$ begins increasing.

At $t = 4$, a second update modifies exception criteria. Exception review load $E(t)$
spikes. The attending physician's available review time falls below sustainable threshold.

At $t = 5$, $D_{II}(t)$ reaches $0.58$, approaching $\theta = 0.60$. Velocity
$\frac{dD_{II}}{dt}$ exceeds $\kappa$. Freeze activates.

Freeze operations suspend further model updates, reduce alert generation to critical-only,
and initiate responsibility handoff to a second reviewer. Recovery rate $r(t)$ increases.
Consumption rate $c(t)$ falls.

By $t = 7$, $B(t)$ has recovered to $0.72$. $D_{II}(t) = 0.28$. The velocity constraint
is restored.

No biological irreversibility occurred. The freeze prevented institutional irreversibility
from collapsing evaluative space before the clinical team could reassert meaningful oversight.

This scenario illustrates that ICG does not regulate clinical outcomes directly. It preserves
the structural conditions under which clinical judgment remains executable.

\section{Relation to Existing Frameworks}

\subsection{Responsibility Frameworks}

Meaningful Human Control specifies structural conditions under which system outputs remain
responsive to relevant human reasons and attributable through intelligible causal chains
(Santoni de Sio \& van den Hoven, 2018).

Reachability-based approaches conceptualize responsibility in terms of avoiding undesirable
global system configurations (Basir, 2026).

A companion framework addressing the temporal fixation of institutional responsibility
under AI-mediated conditions has been developed in parallel work.

ICG operates at a distinct analytical layer. Rather than specifying alignment conditions
or endpoint avoidance, it models the institutional dynamics under which evaluative capacity
may degrade prior to irreversible outcomes.

Where MHC asks whether system outputs remain attributable to human reasons, ICG asks
whether the institutional capacity to evaluate attribution remains intact. These are
complementary questions operating at different temporal layers. MHC addresses the moment
of decision. ICG addresses the structural conditions that must hold before and during
decision-making for MHC conditions to remain satisfiable.

Where reachability frameworks ask whether forbidden configurations remain unreachable,
ICG asks whether the velocity of approach toward such configurations is regulated. An
institution may not have crossed an irreversibility threshold while nonetheless drifting
toward it at a rate that exceeds its corrective capacity.

\subsection{Constraint Logic and Control-Theoretic Analogy}

Control barrier functions preserve admissible state regions by constraining system
trajectories (Ames et al., 2017).

ICG does not extend these frameworks nor operate in physical state spaces. Its object of
analysis is institutional oversight capacity. The analogy is structural: trajectory
constraint replaces outcome maximization.

ICG translates this constraint logic into governance environments where the preserved
object is not mechanical stability but the operational viability of human responsibility.

\subsection{Regulatory and Clinical Governance Context}

The FDA's AI/ML-Based SaMD Action Plan outlines governance mechanisms for adaptive updates,
including predetermined change control plans and post-market monitoring (FDA, 2021a).

Good Machine Learning Practice articulates implementation principles for data management,
validation, and human oversight in medical AI development (FDA, 2021b).

Empirical studies of ICU deployment demonstrate that interpretive and supervisory burdens
materially shape clinical decision environments (Sendak et al., 2020).

These frameworks consistently state that human oversight must be ensured. ICG reformulates
this obligation in dynamic terms: oversight capacity is a state variable subject to
consumption and recovery, and its preservation requires explicit constraints on system
evolution.

\section{Discussion}

Continuous optimization reduces friction. Reduced friction accelerates institutional drift.
When drift exceeds oversight capacity, responsibility remains formally assigned but
substantively hollow.

ICG reintroduces structured friction to preserve evaluative space.

The model is not growth-optimizing but collapse-delaying in the precise sense that it
constrains the velocity of approach toward institutional irreversibility. Its objective
is to preserve the structural conditions under which human responsibility remains real
rather than merely formal.

Existing audit requirements address compliance events after they occur. ICG models oversight
capacity as a dynamic resource whose depletion can precede audit failure. The object of
governance is therefore not the event but the conditions that make events recoverable.

Oversight capacity ultimately reflects an institution's ability to re-evaluate and, if
necessary, replace its own operative assumptions. ICG protects this capacity without
formalizing it as a metaphysical property.

\section{Conclusion: Institutionalized Stopping as Governance Principle}

Irreversibility-Constrained Governance preserves responsibility not by retrospective
attribution but by ex ante trajectory constraint.

\textit{A system that cannot stop cannot remain responsible.}

By constraining the velocity of approach toward institutional irreversibility, governance
prevents responsibility evaporation under continuous revision. Freeze mechanisms are
structured recovery operations that preserve evaluative capacity.

In medical environments, this distinction becomes existential. Biological irreversibility
cannot be undone. Institutional irreversibility must therefore be detected and slowed
before oversight capacity collapses.

This is not a growth-optimization model. It is a collapse-delay model in the precise
sense that it constrains the trajectory toward institutional irreversibility.

Its objective is civilizational: to preserve the structural conditions under which
human beings can still say --- and mean --- ``here we stop.''

\section*{Declarations}

\noindent\textbf{Funding}\quad This research received no external funding.

\noindent\textbf{Conflict of Interest}\quad No conflicts of interest to declare.

\noindent\textbf{Ethical Approval}\quad This manuscript does not involve human subjects
research or identifiable patient data. Institutional review board approval was not required.

\noindent\textbf{Data Availability}\quad No empirical datasets were generated or analyzed
in this study.

\noindent\textbf{AI Use Disclosure}\quad The author used AI-based language assistance tools
for drafting support and structural refinement. All conceptual development, modeling, and
final content decisions were made by the author.

\section*{References}

\begin{description}[leftmargin=1cm, labelindent=0cm, itemsep=4pt]

\item Ames, A.\,D., Coogan, S., Egerstedt, M., Notomista, G., Sreenath, K., \& Tabuada, P.
(2017). Control barrier functions: Theory and applications.
\textit{European Control Conference (ECC)}, 3420--3431.

\item Basir, O. (2026). AI social responsibility as reachability: Execution-level semantics
for the social responsibility stack. \textit{arXiv:2601.02585}.

\item Food and Drug Administration. (2021a). \textit{Artificial intelligence/machine
learning-based software as a medical device action plan.}
U.S. Department of Health and Human Services.

\item Food and Drug Administration. (2021b). \textit{Good machine learning practice for
medical device development: Guiding principles.}
U.S. Department of Health and Human Services.

\item Santoni de Sio, F., \& van den Hoven, J. (2018). Meaningful human control over
autonomous systems: A philosophical account.
\textit{Frontiers in Robotics and AI, 5}, 15.

\item Sendak, M., Elish, M.\,C., Gao, M., Futoma, J., Ratliff, W., Nichols, M.,
Bedoya, A., Balu, S., \& O'Brien, C. (2020).
``The human body is a black box'': Supporting clinical decision-making with deep learning.
\textit{Proceedings of the ACM Conference on Fairness, Accountability, and Transparency},
99--109.

\end{description}

\end{document}
